% 第 30 章骨架：导引框已定，正文待写。写作计划见 docs/PLAN.md“第 30 章”，写法见 docs/STYLE.md。
\chapter{1930 年代的另类出路：纳粹德国、高桥财政与瑞典}\label{ch:alternatives}

\begin{chapterguide}
\guidequestion{纳粹德国、高桥财政下的日本和瑞典都较早走出了萧条，各自靠的是什么？代价记在了哪里？}
\begin{guidegoals}
\item 重建纳粹德国用 MEFO 票据为再军备融资、用外汇管制与 1936 年物价冻结压住公开通胀的机制，并如实说明它服务于侵略战争与掠夺，以及最终的代价；
\item 分析日本 1931 年脱离金本位与日本银行承购国债的效果，以及它的政治结局；
\item 解释瑞典 1931 年后的价格水平目标；
\item 与美国新政和苏联五年计划对照，讨论“政策实践先于理论”的含义。
\end{guidegoals}
\guideprereq{\chref{ch:depressionworld}\enspace\nameref{ch:depressionworld}}
\end{chapterguide}

\chapterstub
